\documentclass[border=8pt,tikz]{standalone}
\usepackage[T1]{fontenc}
\usepackage[utf8]{inputenc}
\usepackage[portuguese]{babel}
\usepackage{lmodern}
\usepackage{amsmath}
\usepackage{tikz}
\usetikzlibrary{arrows.meta,positioning,shapes.geometric,calc,fit,backgrounds,decorations.pathreplacing}

\definecolor{verdeEsc}{HTML}{1B5E20}
\definecolor{verde}{HTML}{2E7D32}
\definecolor{verdeClaro}{HTML}{EDF5EE}
\definecolor{verdeMed}{HTML}{CFE6D2}
\definecolor{azul}{HTML}{1565C0}
\definecolor{azulClaro}{HTML}{E3EEFA}
\definecolor{laranja}{HTML}{E65100}
\definecolor{laranjaClaro}{HTML}{FBE9D8}
\definecolor{ambar}{HTML}{B07300}
\definecolor{ambarClaro}{HTML}{FBF3DA}
\definecolor{vermelho}{HTML}{B3261E}
\definecolor{vermelhoClaro}{HTML}{F8E2E0}
\definecolor{cinza}{HTML}{4A5560}

\tikzset{
  caixa/.style={draw=verde, fill=verdeClaro, rounded corners=1.5pt, line width=.6pt,
                align=center, inner sep=4pt, font=\small, minimum height=9mm},
  destaque/.style={caixa, fill=verdeMed, draw=verdeEsc, line width=1pt},
  raiz/.style={caixa, fill=verde, draw=verdeEsc, text=white, line width=1pt, font=\small\bfseries},
  decisao/.style={draw=ambar, fill=ambarClaro, diamond, aspect=1.7, align=center,
                  inner sep=1pt, font=\footnotesize, line width=.6pt},
  azulbox/.style={caixa, fill=azulClaro, draw=azul},
  laranjabox/.style={caixa, fill=laranjaClaro, draw=laranja},
  vermelhobox/.style={caixa, fill=vermelhoClaro, draw=vermelho},
  seta/.style={-{Latex[length=2.4mm,width=1.8mm]}, draw=cinza, line width=.6pt},
  setaverde/.style={seta, draw=verdeEsc},
  setavermelha/.style={seta, draw=vermelho},
  rot/.style={font=\scriptsize, inner sep=1.5pt, fill=white},
  grupo/.style={draw=cinza, dashed, rounded corners=3pt, line width=.5pt, inner sep=7pt},
  titulogrupo/.style={font=\footnotesize\bfseries, text=cinza}
}
\hyphenpenalty=10000
\exhyphenpenalty=10000
\tolerance=9999
\emergencystretch=3em

\begin{document}

\begin{tikzpicture}[node distance=8mm]
  \node[raiz, text width=34mm] (S) {Feição erosiva\\ identificada};
  \node[decisao, below=10mm of S, text width=18mm] (Q1) {Profundidade\\ $>$ 1,5 m?};

  \node[decisao, below left=12mm and 30mm of Q1, text width=16mm] (Q2) {Presença\\ de água?};
  \node[decisao, below right=12mm and 30mm of Q1, text width=16mm] (Q3) {Inclinação\\ $>$ 45\textdegree?};

  \node[destaque, text width=28mm, below left=13mm and 6mm of Q2] (T1) {Gabião vivo\\ e geodreno};
  \node[destaque, text width=28mm, below right=13mm and 6mm of Q2] (T2) {Parede Krainer\\ e revegetação};

  \node[decisao, below left=13mm and 2mm of Q3, text width=16mm] (Q4) {Solo\\ coesivo?};
  \node[destaque, text width=28mm, below right=13mm and 6mm of Q3] (T3) {Paliçadas\\ e biomanta};

  \node[destaque, text width=28mm, below left=13mm and 4mm of Q4] (T4) {Canaleta verde\\ e hidrossemeadura};
  \node[destaque, text width=28mm, below right=13mm and 4mm of Q4] (T5) {Enrocamento\\ vegetado};

  \draw[seta] (S) -- (Q1);
  \draw[seta] (Q1) -- node[rot, above left] {sim} (Q2);
  \draw[seta] (Q1) -- node[rot, above right] {não} (Q3);
  \draw[seta] (Q2) -- node[rot, left] {sim} (T1);
  \draw[seta] (Q2) -- node[rot, right] {não} (T2);
  \draw[seta] (Q3) -- node[rot, left] {sim} (Q4);
  \draw[seta] (Q3) -- node[rot, right] {não} (T3);
  \draw[seta] (Q4) -- node[rot, left] {sim} (T4);
  \draw[seta] (Q4) -- node[rot, right] {não} (T5);
\end{tikzpicture}

\end{document}
